\documentclass[a4paper,11pt]{ltjsarticle}
\usepackage[haranoaji]{luatexja-preset}
\usepackage{amsmath,amssymb}
\usepackage[top=12mm,bottom=10mm,left=14mm,right=14mm]{geometry}
\usepackage{array}
\usepackage{tikz}
\usetikzlibrary{calc}
\pagestyle{empty}
\setlength{\parindent}{0pt}
\newlength{\qcol}\setlength{\qcol}{0.47\textwidth}
\newlength{\qgap}
\newlength{\anscolw}\setlength{\anscolw}{0.30\textwidth}
\newlength{\ansh}
\newcommand{\Q}[2]{\parbox[t]{\qcol}{\raggedright (#1)\hspace{0.6em}$\displaystyle #2$}}
\newcommand{\QR}[4]{\Q{#1}{#2}\hfill\Q{#3}{#4}\par\vspace{\qgap}}
\newcommand{\anscell}[1]{\rule[-\ansdrop]{0pt}{\ansh}(#1)}
\newlength{\ansdrop}

\begin{document}
{\large\bfseries 計算問題　18}\hfill 名前(\underline{\hspace{32mm}})\par\vspace{3mm}
■（1）〜（16）の計算をしなさい。（17）〜（21）は方程式を解きなさい。\par\vspace{3mm}
\setlength{\qgap}{5.4mm}
\QR{1}{108x\div(-9)}{2}{(-8)\times14a}
\QR{3}{-13+(-8)}{4}{(-56)\div(+4)}
\QR{5}{(5-2x)+(3x-6)}{6}{(-6)-(-9)+(-5)}
\QR{7}{8-(-3)\div\left(-\dfrac{3}{5}\right)}{8}{6\div5\times(-10)}
\QR{9}{(-14)\times(-16)}{10}{(15a-24)\div3}
\QR{11}{\dfrac{3}{5}\times\left\{-\dfrac{1}{2}+\left(-\dfrac{7}{6}\right)\right\}}{12}{\dfrac{2}{5}-\dfrac{3}{2}+(-0.7)}
\QR{13}{(-5)^2\times(-1)^3}{14}{\dfrac{1}{2}(8x-4)-\dfrac{1}{3}(9x-3)}
\QR{15}{3\times(-9)-16\div(-2)}{16}{(-29)-(-14)}
\QR{17}{0.7x+0.1=-2}{18}{\dfrac{1}{2}x+5=\dfrac{1}{6}x+3}
\QR{19}{4(x-6)=3(x-1)}{20}{-3x+7=-5}
\vspace{1mm}
\Q{21}{\dfrac{x+2}{3}=\dfrac{x-2}{2}}\par\vspace{3mm}
\setlength{\ansh}{9.5mm}\setlength{\ansdrop}{6mm}%
\begin{center}\begin{tabular}{|p{\anscolw}|p{\anscolw}|p{\anscolw}|}\hline
\anscell{1} & \anscell{2} & \anscell{3} \\\hline
\anscell{4} & \anscell{5} & \anscell{6} \\\hline
\anscell{7} & \anscell{8} & \anscell{9} \\\hline
\anscell{10} & \anscell{11} & \anscell{12} \\\hline
\anscell{13} & \anscell{14} & \anscell{15} \\\hline
\anscell{16} & \anscell{17} & \anscell{18} \\\hline
\anscell{19} & \anscell{20} & \anscell{21} \\\hline
\end{tabular}\end{center}

\end{document}
